\documentclass{article}
\usepackage{/globals/de}
\usepackage{/globals/anki}
\ankiprefix{begriff-staat}
\ankitopic{Begriff Staat}
\ankishow
\begin{document}
\section{Begriff: Staat} 
Ein Staat ist ein komplexes Konstrukt, kann aber anhand von ein paar Merkmalen erkannt werden. Ein Staat hat innerhalb eines \emph{Staatsgebiet}es die \emph{Staatsgewalt} über das dort lebende \emph{Staatsvolk}.

Mit der Staatsgewalt hat der Staat das Monopol über die Gesetzgebung und die Gewalt der Exekutive.
\ankinote{def}{Definition}{Staatsgebiet, -gewalt, -volk}

\subsection{Prinzipien}
\begin{description}
 \item[Demokratie] Der Staat sollte durch demokratische Wahlen legitimiert sein.
 \item[Rechtsstaatlichkeit] Der Staat sollte nicht willkürlich sein, sollte seine eigenen Regeln befolgen.
 \item[Gewaltmonopol] Der Staat hat das Monopol der Gewaltanwendung.
 \item[Unabhängige Verwaltung] Bürokratie, \zb Steuern, sind für die Existenz notwendig.
\end{description}
\ankinote{prinzipien}{Prinzipien}{Demokratie, Rechtsstaatlichkeit, Gewaltmonopol, Unabhängige Verwaltung}

\subsection{Staatsformen}
\begin{description}
 \item[Demokratien] Es gibt demokratische Wahlen, Beteiligung.
 \item[Autoritäre Staaten] Eine Person oder Personengruppe hat die Macht, es gibt nur wenig Beteiligung.
 \item[Diktaturen] Eine Person oder Familie hat die Macht, Beteiligung wird unterdrückt, es gibt keine fairen Wahlen.
\end{description}
\ankinote{formen}{Staatsformen}{Demokratie, Autoritär, Diktaturen (unterdrücking)}
\end{document}